% BOOK-FILE-ID: f-e862ee83e454
% subsubsection: 逆数型  \, $\displaystyle a_{n+1}=\frac{pa_n}{qa_n+r}$

\subsubsection{\texorpdfstring{逆数型   $a_{n+1}=\frac{pa_n}{qa_n+r}$}{逆数型}}\label{節：逆数型}
おそらくよく扱われる漸化式の中で一番むずかしいと呼ばれるであろう分数型の内の1パターン目を紹介する.
一般的な形は
\[
  a_{n+1}=\frac{pa_n}{qa_n+r}
\]
で与えられる.

以下では,\,$p\ne0$かつ$a_n\ne0$であり,\,$qa_n+r\ne0$として議論する.\footnote{$a_1=0$は逆数を取る前に分ける.\,$r\ne0$ならすべての項が$0$となる.\,$r=0$なら$a_2$を求める段階で分母が$0$となり,そのままでは定義できない.}

少々発想しにくい式変形であるため,結果からお見せしよう.
この変形の肝は逆数をとることである.
両辺の逆数を取り,\,$b_n=\frac{1}{a_n}$と置くと,
\begin{align*}
  \frac{1}{a_{n+1}}&=\frac{qa_n+r}{pa_n}
  &=\frac{q}{p}+\frac{r}{p}\cdot\frac{1}{a_n},\\
  b_{n+1}&=\frac{r}{p}b_n+\frac{q}{p}
\end{align*}
ここから特性方程式型または等差型として$b_n$を求め,最後に$a_n=\frac{1}{b_n}$へ戻す.
ただし$b_n=0$となる項は逆数を戻せないため,元の漸化式が定義されるかも確認する.

\begin{bookpoint}
  \textbf{POINT}\par
  分数の両方に$a_n$が存在し,分子側が$a_n$の定数倍のみであるときには逆数を取って約分
\end{bookpoint}

なお,
\[
  a_{n+1}=\frac{pa_n+q}{ra_n}
\]
の場合は\booknavref{節：一次分数型}{一次分数型}を参照すること.

\noindent \bookblocklink{example-reciprocal-1}{problem}{answer}{【例題】}\par\nobreak
\begin{enumerate}[label=(\arabic*),parsep=3pt,format=\bookitemlink{example-reciprocal-1}{problem}{answer}]
  \item $a_1=1,\, a_{n+1}=\frac{a_n}{a_n+1}$
  \item $a_1=1,\, a_{n+1}=\frac{-a_n}{2a_n+1}$
\end{enumerate}

\noindent\bookblocklink{example-reciprocal-1}{hint}{problem}{【方針】}\par\nobreak
\begin{enumerate}[label=(\arabic*),parsep=3pt,format=\bookitemlink{example-reciprocal-1}{hint}{problem}]
  \item 逆数を取ると,\,$\frac{1}{a_n}$が毎回$1$ずつ増える形になります.
  \item 逆数を取って一次の\bookterm{recurrence}{漸化式}に帰着させましょう.
\end{enumerate}

\noindent\bookblocklink{example-reciprocal-1}{answer}{problem}{【解答】}\par\nobreak
\begin{enumerate}[label=(\arabic*),parsep=3pt,format=\bookitemlink{example-reciprocal-1}{answer}{problem}]
  \item $a_n=\frac{1}{n}$
  \item $a_n=\frac{1}{2(-1)^{n-1}-1}$
\end{enumerate}

\noindent\bookblocklink{example-reciprocal-1}{solution}{problem}{【解法】}\par\nobreak
\begin{enumerate}[label=(\arabic*),parsep=3pt,format=\bookitemlink{example-reciprocal-1}{solution}{problem}]
  \item $a_1=1,\, a_{n+1}=\frac{a_n}{a_n+1}$\\
    正の項と分母が続くので,逆数を取れる.\,$b_n=\frac{1}{a_n}$と置くと,
    \begin{align*}
      &\frac{1}{a_{n+1}}=1+\frac{1}{a_n},\\
      &b_{n+1}=b_n+1,\quad b_1=1\\
      &\therefore \quad b_n=n,\quad a_n=\frac1n
    \end{align*}
  \item $a_1=1,\, a_{n+1}=\frac{-a_n}{2a_n+1}$\\
    初めの数項は$1\longmapsto-\frac{1}{3}\longmapsto1$となる.
    項も分母も零でないので,\,$b_n=\frac{1}{a_n}$と置くと,
    \begin{align*}
      &\frac{1}{a_{n+1}}=-2-\frac{1}{a_n},\\
      &b_{n+1}=-b_n-2,\quad b_1=1
    \end{align*}
    \bookterm{fixed}{不動点}は$-1$なので,\,$c_n=b_n+1$と置き,定数項を消す.
    \begin{align*}
      &c_{n+1}=-c_n,\quad c_1=2\\
      &\therefore \quad b_n=c_n-1=2(-1)^{n-1}-1,\\
      &a_n=\frac{1}{2(-1)^{n-1}-1}
    \end{align*}
\end{enumerate}
